% =======================================================
% CONFIGURAÇÕES E IMPRESSÃO DAS LISTAS (TUDO EM UM ARQUIVO)
% Pode ser chamado com % =======================================================
% CONFIGURAÇÕES E IMPRESSÃO DAS LISTAS (TUDO EM UM ARQUIVO)
% Pode ser chamado com % =======================================================
% CONFIGURAÇÕES E IMPRESSÃO DAS LISTAS (TUDO EM UM ARQUIVO)
% Pode ser chamado com % =======================================================
% CONFIGURAÇÕES E IMPRESSÃO DAS LISTAS (TUDO EM UM ARQUIVO)
% Pode ser chamado com \input{listas} APÓS o \begin{document}
	% =======================================================
	
	\makeatletter
	
	%------------------------------------------------------------
	% 1. AJUSTES DO SUMÁRIO
	%------------------------------------------------------------
	\renewcommand{\cftsubsecfont}{\itshape\bfseries}
	\renewcommand{\cftsubsecpagefont}{\normalfont}
	
	%------------------------------------------------------------
	% 2. AJUSTES DOS PREFIXOS DAS LISTAS
	%------------------------------------------------------------
	\@ifundefined{cftfigpresnum}{}{
		\renewcommand{\cftfigpresnum}{\figurename~}
		\renewcommand{\cftfigaftersnum}{~--~}
		\settowidth{\cftfignumwidth}{\cftfigpresnum 000~--~}
	}
	\@ifundefined{cfttabpresnum}{}{
		\renewcommand{\cfttabpresnum}{\tablename~}
		\renewcommand{\cfttabaftersnum}{~--~}
		\settowidth{\cfttabnumwidth}{\cfttabpresnum 000~--~}
	}
	\@ifundefined{cftquadropresnum}{}{
		\renewcommand{\cftquadropresnum}{Quadro~}
		\renewcommand{\cftquadroaftersnum}{~--~}
		\settowidth{\cftquadronumwidth}{\cftquadropresnum 000~--~}
	}
	\@ifundefined{cftgraficopresnum}{}{
		\renewcommand{\cftgraficopresnum}{Gráfico~}
		\renewcommand{\cftgraficoaftersnum}{~--~}
		\settowidth{\cftgraficonumwidth}{\cftgraficopresnum 000~--~}
	}
	
	%------------------------------------------------------------
	% 3. REGISTRO E DETECÇÃO AUTOMÁTICA DE ELEMENTOS NO TEXTO
	%------------------------------------------------------------
	\let\native@addcontentsline\addcontentsline
	\renewcommand{\addcontentsline}[3]{%
		\immediate\write\@mainaux{\string\gdef\string\temitem@#1{1}}%
		\native@addcontentsline{#1}{#2}{#3}%
	}
	
	% Travas lógicas para escrever no .aux apenas 1 vez e evitar loops
	\def\temitem@acr@written{0}
	\newcommand{\RegistrarUsoAcr}{%
		\ifnum\temitem@acr@written=0
		\immediate\write\@mainaux{\string\gdef\string\temitem@acr{1}}%
		\global\def\temitem@acr@written{1}%
		\fi
	}
	
	\def\temitem@nls@written{0}
	\newcommand{\RegistrarUsoNls}{%
		\ifnum\temitem@nls@written=0
		\immediate\write\@mainaux{\string\gdef\string\temitem@nls{1}}%
		\global\def\temitem@nls@written{1}%
		\fi
	}
	
	% Monitora os comandos padrões do glossaries e o SEU comando customizado \sigla
	\@ifpackageloaded{glossaries}{
		\AddToHook{cmd/glsadd/before}{\RegistrarUsoAcr}
		\AddToHook{cmd/gls/before}{\RegistrarUsoAcr}
	}{}
	% Gatilho crucial para quem usa o comando \sigla{}{}:
	\AddToHook{cmd/sigla/before}{\RegistrarUsoAcr}
	
	\@ifpackageloaded{nomencl}{
		\AddToHook{cmd/nomenclature/before}{\RegistrarUsoNls}
	}{}
	
	\newcommand{\SeTemItem}[2]{%
		\def\elemento@achado{0}%
		\@ifundefined{temitem@#1}{}{%
			\def\elemento@achado{1}%
		}%
		\@ifundefined{ext@#1}{}{%
			\@ifundefined{temitem@\csname ext@#1\endcsname}{}{%
				\def\elemento@achado{1}%
			}%
		}%
		\ifnum\elemento@achado=1\relax
		#2%
		\fi
	}
	
	%------------------------------------------------------------
	% 4. LISTAS SEM TÍTULO AUTOMÁTICO E FORMATAÇÃO
	%------------------------------------------------------------
	\renewcommand{\listoffigures}{\@starttoc{lof}}
	\renewcommand{\listoftables}{\@starttoc{lot}}
	\renewcommand{\listof}[2]{\@starttoc{\csname ext@#1\endcsname}}
	
	\newcommand{\PreTextualTitulo}[1]{%
		\begin{center}
			\bfseries\MakeUppercase{#1}
		\end{center}
	}
	
	\makeatother
	
	% =======================================================
	% 5. IMPRESSÃO DAS LISTAS E SUMÁRIO
	% =======================================================
	
	% LISTA DE ILUSTRAÇÕES
	\SeTemItem{lof}{%
		\PreTextualTitulo{Lista de Ilustrações}
		\listoffigures
		\thispagestyle{empty}
		\newpage
	}
	
	% LISTA DE TABELAS
	\SeTemItem{lot}{%
		\PreTextualTitulo{Lista de Tabelas}
		\listoftables
		\thispagestyle{empty}
		\newpage
	}
	
	% LISTA DE ABREVIATURAS E SIGLAS
	\SeTemItem{acr}{%
		% Trava de segurança extra para garantir que não haja registros durante a impressão
		\global\def\temitem@acr@written{1}
		\PreTextualTitulo{Lista de Abreviaturas e Siglas}
		\printglossary[
		type=\acronymtype,
		style=siglasABNT,
		title={}
		]
		\thispagestyle{empty}
		\newpage
	}
	
	% LISTA DE SÍMBOLOS
	\SeTemItem{nls}{%
		\global\def\temitem@nls@written{1}
		\printnomenclature
		\thispagestyle{empty}
		\newpage
	}
	
	% LISTA DE QUADROS
	\SeTemItem{quadro}{%
		\PreTextualTitulo{Lista de Quadros}
		\listof{quadro}{}
		\thispagestyle{empty}
		\newpage
	}
	
	% LISTA DE GRÁFICOS
	\SeTemItem{grafico}{%
		\PreTextualTitulo{Lista de Gráficos}
		\listof{grafico}{}
		\thispagestyle{empty}
		\newpage
	}
	
	% SUMÁRIO
	\tableofcontents
	\thispagestyle{empty}
	
	% =======================================================
	% 6. CORREÇÃO DOS CONTADORES (ABNT)
	% =======================================================
	\setcounter{table}{0}
	\setcounter{figure}{0}
	\makeatletter
	\@ifundefined{c@quadro}{}{\setcounter{quadro}{0}}
	\@ifundefined{c@grafico}{}{\setcounter{grafico}{0}}
	\makeatother APÓS o \begin{document}
	% =======================================================
	
	\makeatletter
	
	%------------------------------------------------------------
	% 1. AJUSTES DO SUMÁRIO
	%------------------------------------------------------------
	\renewcommand{\cftsubsecfont}{\itshape\bfseries}
	\renewcommand{\cftsubsecpagefont}{\normalfont}
	
	%------------------------------------------------------------
	% 2. AJUSTES DOS PREFIXOS DAS LISTAS
	%------------------------------------------------------------
	\@ifundefined{cftfigpresnum}{}{
		\renewcommand{\cftfigpresnum}{\figurename~}
		\renewcommand{\cftfigaftersnum}{~--~}
		\settowidth{\cftfignumwidth}{\cftfigpresnum 000~--~}
	}
	\@ifundefined{cfttabpresnum}{}{
		\renewcommand{\cfttabpresnum}{\tablename~}
		\renewcommand{\cfttabaftersnum}{~--~}
		\settowidth{\cfttabnumwidth}{\cfttabpresnum 000~--~}
	}
	\@ifundefined{cftquadropresnum}{}{
		\renewcommand{\cftquadropresnum}{Quadro~}
		\renewcommand{\cftquadroaftersnum}{~--~}
		\settowidth{\cftquadronumwidth}{\cftquadropresnum 000~--~}
	}
	\@ifundefined{cftgraficopresnum}{}{
		\renewcommand{\cftgraficopresnum}{Gráfico~}
		\renewcommand{\cftgraficoaftersnum}{~--~}
		\settowidth{\cftgraficonumwidth}{\cftgraficopresnum 000~--~}
	}
	
	%------------------------------------------------------------
	% 3. REGISTRO E DETECÇÃO AUTOMÁTICA DE ELEMENTOS NO TEXTO
	%------------------------------------------------------------
	\let\native@addcontentsline\addcontentsline
	\renewcommand{\addcontentsline}[3]{%
		\immediate\write\@mainaux{\string\gdef\string\temitem@#1{1}}%
		\native@addcontentsline{#1}{#2}{#3}%
	}
	
	% Travas lógicas para escrever no .aux apenas 1 vez e evitar loops
	\def\temitem@acr@written{0}
	\newcommand{\RegistrarUsoAcr}{%
		\ifnum\temitem@acr@written=0
		\immediate\write\@mainaux{\string\gdef\string\temitem@acr{1}}%
		\global\def\temitem@acr@written{1}%
		\fi
	}
	
	\def\temitem@nls@written{0}
	\newcommand{\RegistrarUsoNls}{%
		\ifnum\temitem@nls@written=0
		\immediate\write\@mainaux{\string\gdef\string\temitem@nls{1}}%
		\global\def\temitem@nls@written{1}%
		\fi
	}
	
	% Monitora os comandos padrões do glossaries e o SEU comando customizado \sigla
	\@ifpackageloaded{glossaries}{
		\AddToHook{cmd/glsadd/before}{\RegistrarUsoAcr}
		\AddToHook{cmd/gls/before}{\RegistrarUsoAcr}
	}{}
	% Gatilho crucial para quem usa o comando \sigla{}{}:
	\AddToHook{cmd/sigla/before}{\RegistrarUsoAcr}
	
	\@ifpackageloaded{nomencl}{
		\AddToHook{cmd/nomenclature/before}{\RegistrarUsoNls}
	}{}
	
	\newcommand{\SeTemItem}[2]{%
		\def\elemento@achado{0}%
		\@ifundefined{temitem@#1}{}{%
			\def\elemento@achado{1}%
		}%
		\@ifundefined{ext@#1}{}{%
			\@ifundefined{temitem@\csname ext@#1\endcsname}{}{%
				\def\elemento@achado{1}%
			}%
		}%
		\ifnum\elemento@achado=1\relax
		#2%
		\fi
	}
	
	%------------------------------------------------------------
	% 4. LISTAS SEM TÍTULO AUTOMÁTICO E FORMATAÇÃO
	%------------------------------------------------------------
	\renewcommand{\listoffigures}{\@starttoc{lof}}
	\renewcommand{\listoftables}{\@starttoc{lot}}
	\renewcommand{\listof}[2]{\@starttoc{\csname ext@#1\endcsname}}
	
	\newcommand{\PreTextualTitulo}[1]{%
		\begin{center}
			\bfseries\MakeUppercase{#1}
		\end{center}
	}
	
	\makeatother
	
	% =======================================================
	% 5. IMPRESSÃO DAS LISTAS E SUMÁRIO
	% =======================================================
	
	% LISTA DE ILUSTRAÇÕES
	\SeTemItem{lof}{%
		\PreTextualTitulo{Lista de Ilustrações}
		\listoffigures
		\thispagestyle{empty}
		\newpage
	}
	
	% LISTA DE TABELAS
	\SeTemItem{lot}{%
		\PreTextualTitulo{Lista de Tabelas}
		\listoftables
		\thispagestyle{empty}
		\newpage
	}
	
	% LISTA DE ABREVIATURAS E SIGLAS
	\SeTemItem{acr}{%
		% Trava de segurança extra para garantir que não haja registros durante a impressão
		\global\def\temitem@acr@written{1}
		\PreTextualTitulo{Lista de Abreviaturas e Siglas}
		\printglossary[
		type=\acronymtype,
		style=siglasABNT,
		title={}
		]
		\thispagestyle{empty}
		\newpage
	}
	
	% LISTA DE SÍMBOLOS
	\SeTemItem{nls}{%
		\global\def\temitem@nls@written{1}
		\printnomenclature
		\thispagestyle{empty}
		\newpage
	}
	
	% LISTA DE QUADROS
	\SeTemItem{quadro}{%
		\PreTextualTitulo{Lista de Quadros}
		\listof{quadro}{}
		\thispagestyle{empty}
		\newpage
	}
	
	% LISTA DE GRÁFICOS
	\SeTemItem{grafico}{%
		\PreTextualTitulo{Lista de Gráficos}
		\listof{grafico}{}
		\thispagestyle{empty}
		\newpage
	}
	
	% SUMÁRIO
	\tableofcontents
	\thispagestyle{empty}
	
	% =======================================================
	% 6. CORREÇÃO DOS CONTADORES (ABNT)
	% =======================================================
	\setcounter{table}{0}
	\setcounter{figure}{0}
	\makeatletter
	\@ifundefined{c@quadro}{}{\setcounter{quadro}{0}}
	\@ifundefined{c@grafico}{}{\setcounter{grafico}{0}}
	\makeatother APÓS o \begin{document}
	% =======================================================
	
	\makeatletter
	
	%------------------------------------------------------------
	% 1. AJUSTES DO SUMÁRIO
	%------------------------------------------------------------
	\renewcommand{\cftsubsecfont}{\itshape\bfseries}
	\renewcommand{\cftsubsecpagefont}{\normalfont}
	
	%------------------------------------------------------------
	% 2. AJUSTES DOS PREFIXOS DAS LISTAS
	%------------------------------------------------------------
	\@ifundefined{cftfigpresnum}{}{
		\renewcommand{\cftfigpresnum}{\figurename~}
		\renewcommand{\cftfigaftersnum}{~--~}
		\settowidth{\cftfignumwidth}{\cftfigpresnum 000~--~}
	}
	\@ifundefined{cfttabpresnum}{}{
		\renewcommand{\cfttabpresnum}{\tablename~}
		\renewcommand{\cfttabaftersnum}{~--~}
		\settowidth{\cfttabnumwidth}{\cfttabpresnum 000~--~}
	}
	\@ifundefined{cftquadropresnum}{}{
		\renewcommand{\cftquadropresnum}{Quadro~}
		\renewcommand{\cftquadroaftersnum}{~--~}
		\settowidth{\cftquadronumwidth}{\cftquadropresnum 000~--~}
	}
	\@ifundefined{cftgraficopresnum}{}{
		\renewcommand{\cftgraficopresnum}{Gráfico~}
		\renewcommand{\cftgraficoaftersnum}{~--~}
		\settowidth{\cftgraficonumwidth}{\cftgraficopresnum 000~--~}
	}
	
	%------------------------------------------------------------
	% 3. REGISTRO E DETECÇÃO AUTOMÁTICA DE ELEMENTOS NO TEXTO
	%------------------------------------------------------------
	\let\native@addcontentsline\addcontentsline
	\renewcommand{\addcontentsline}[3]{%
		\immediate\write\@mainaux{\string\gdef\string\temitem@#1{1}}%
		\native@addcontentsline{#1}{#2}{#3}%
	}
	
	% Travas lógicas para escrever no .aux apenas 1 vez e evitar loops
	\def\temitem@acr@written{0}
	\newcommand{\RegistrarUsoAcr}{%
		\ifnum\temitem@acr@written=0
		\immediate\write\@mainaux{\string\gdef\string\temitem@acr{1}}%
		\global\def\temitem@acr@written{1}%
		\fi
	}
	
	\def\temitem@nls@written{0}
	\newcommand{\RegistrarUsoNls}{%
		\ifnum\temitem@nls@written=0
		\immediate\write\@mainaux{\string\gdef\string\temitem@nls{1}}%
		\global\def\temitem@nls@written{1}%
		\fi
	}
	
	% Monitora os comandos padrões do glossaries e o SEU comando customizado \sigla
	\@ifpackageloaded{glossaries}{
		\AddToHook{cmd/glsadd/before}{\RegistrarUsoAcr}
		\AddToHook{cmd/gls/before}{\RegistrarUsoAcr}
	}{}
	% Gatilho crucial para quem usa o comando \sigla{}{}:
	\AddToHook{cmd/sigla/before}{\RegistrarUsoAcr}
	
	\@ifpackageloaded{nomencl}{
		\AddToHook{cmd/nomenclature/before}{\RegistrarUsoNls}
	}{}
	
	\newcommand{\SeTemItem}[2]{%
		\def\elemento@achado{0}%
		\@ifundefined{temitem@#1}{}{%
			\def\elemento@achado{1}%
		}%
		\@ifundefined{ext@#1}{}{%
			\@ifundefined{temitem@\csname ext@#1\endcsname}{}{%
				\def\elemento@achado{1}%
			}%
		}%
		\ifnum\elemento@achado=1\relax
		#2%
		\fi
	}
	
	%------------------------------------------------------------
	% 4. LISTAS SEM TÍTULO AUTOMÁTICO E FORMATAÇÃO
	%------------------------------------------------------------
	\renewcommand{\listoffigures}{\@starttoc{lof}}
	\renewcommand{\listoftables}{\@starttoc{lot}}
	\renewcommand{\listof}[2]{\@starttoc{\csname ext@#1\endcsname}}
	
	\newcommand{\PreTextualTitulo}[1]{%
		\begin{center}
			\bfseries\MakeUppercase{#1}
		\end{center}
	}
	
	\makeatother
	
	% =======================================================
	% 5. IMPRESSÃO DAS LISTAS E SUMÁRIO
	% =======================================================
	
	% LISTA DE ILUSTRAÇÕES
	\SeTemItem{lof}{%
		\PreTextualTitulo{Lista de Ilustrações}
		\listoffigures
		\thispagestyle{empty}
		\newpage
	}
	
	% LISTA DE TABELAS
	\SeTemItem{lot}{%
		\PreTextualTitulo{Lista de Tabelas}
		\listoftables
		\thispagestyle{empty}
		\newpage
	}
	
	% LISTA DE ABREVIATURAS E SIGLAS
	\SeTemItem{acr}{%
		% Trava de segurança extra para garantir que não haja registros durante a impressão
		\global\def\temitem@acr@written{1}
		\PreTextualTitulo{Lista de Abreviaturas e Siglas}
		\printglossary[
		type=\acronymtype,
		style=siglasABNT,
		title={}
		]
		\thispagestyle{empty}
		\newpage
	}
	
	% LISTA DE SÍMBOLOS
	\SeTemItem{nls}{%
		\global\def\temitem@nls@written{1}
		\printnomenclature
		\thispagestyle{empty}
		\newpage
	}
	
	% LISTA DE QUADROS
	\SeTemItem{quadro}{%
		\PreTextualTitulo{Lista de Quadros}
		\listof{quadro}{}
		\thispagestyle{empty}
		\newpage
	}
	
	% LISTA DE GRÁFICOS
	\SeTemItem{grafico}{%
		\PreTextualTitulo{Lista de Gráficos}
		\listof{grafico}{}
		\thispagestyle{empty}
		\newpage
	}
	
	% SUMÁRIO
	\tableofcontents
	\thispagestyle{empty}
	
	% =======================================================
	% 6. CORREÇÃO DOS CONTADORES (ABNT)
	% =======================================================
	\setcounter{table}{0}
	\setcounter{figure}{0}
	\makeatletter
	\@ifundefined{c@quadro}{}{\setcounter{quadro}{0}}
	\@ifundefined{c@grafico}{}{\setcounter{grafico}{0}}
	\makeatother APÓS o \begin{document}
	% =======================================================
	
	\makeatletter
	
	%------------------------------------------------------------
	% 1. AJUSTES DO SUMÁRIO
	%------------------------------------------------------------
	\renewcommand{\cftsubsecfont}{\itshape\bfseries}
	\renewcommand{\cftsubsecpagefont}{\normalfont}
	
	%------------------------------------------------------------
	% 2. AJUSTES DOS PREFIXOS DAS LISTAS
	%------------------------------------------------------------
	\@ifundefined{cftfigpresnum}{}{
		\renewcommand{\cftfigpresnum}{\figurename~}
		\renewcommand{\cftfigaftersnum}{~--~}
		\settowidth{\cftfignumwidth}{\cftfigpresnum 000~--~}
	}
	\@ifundefined{cfttabpresnum}{}{
		\renewcommand{\cfttabpresnum}{\tablename~}
		\renewcommand{\cfttabaftersnum}{~--~}
		\settowidth{\cfttabnumwidth}{\cfttabpresnum 000~--~}
	}
	\@ifundefined{cftquadropresnum}{}{
		\renewcommand{\cftquadropresnum}{Quadro~}
		\renewcommand{\cftquadroaftersnum}{~--~}
		\settowidth{\cftquadronumwidth}{\cftquadropresnum 000~--~}
	}
	\@ifundefined{cftgraficopresnum}{}{
		\renewcommand{\cftgraficopresnum}{Gráfico~}
		\renewcommand{\cftgraficoaftersnum}{~--~}
		\settowidth{\cftgraficonumwidth}{\cftgraficopresnum 000~--~}
	}
	
	%------------------------------------------------------------
	% 3. REGISTRO E DETECÇÃO AUTOMÁTICA DE ELEMENTOS NO TEXTO
	%------------------------------------------------------------
	\let\native@addcontentsline\addcontentsline
	\renewcommand{\addcontentsline}[3]{%
		\immediate\write\@mainaux{\string\gdef\string\temitem@#1{1}}%
		\native@addcontentsline{#1}{#2}{#3}%
	}
	
	% Travas lógicas para escrever no .aux apenas 1 vez e evitar loops
	\def\temitem@acr@written{0}
	\newcommand{\RegistrarUsoAcr}{%
		\ifnum\temitem@acr@written=0
		\immediate\write\@mainaux{\string\gdef\string\temitem@acr{1}}%
		\global\def\temitem@acr@written{1}%
		\fi
	}
	
	\def\temitem@nls@written{0}
	\newcommand{\RegistrarUsoNls}{%
		\ifnum\temitem@nls@written=0
		\immediate\write\@mainaux{\string\gdef\string\temitem@nls{1}}%
		\global\def\temitem@nls@written{1}%
		\fi
	}
	
	% Monitora os comandos padrões do glossaries e o SEU comando customizado \sigla
	\@ifpackageloaded{glossaries}{
		\AddToHook{cmd/glsadd/before}{\RegistrarUsoAcr}
		\AddToHook{cmd/gls/before}{\RegistrarUsoAcr}
	}{}
	% Gatilho crucial para quem usa o comando \sigla{}{}:
	\AddToHook{cmd/sigla/before}{\RegistrarUsoAcr}
	
	\@ifpackageloaded{nomencl}{
		\AddToHook{cmd/nomenclature/before}{\RegistrarUsoNls}
	}{}
	
	\newcommand{\SeTemItem}[2]{%
		\def\elemento@achado{0}%
		\@ifundefined{temitem@#1}{}{%
			\def\elemento@achado{1}%
		}%
		\@ifundefined{ext@#1}{}{%
			\@ifundefined{temitem@\csname ext@#1\endcsname}{}{%
				\def\elemento@achado{1}%
			}%
		}%
		\ifnum\elemento@achado=1\relax
		#2%
		\fi
	}
	
	%------------------------------------------------------------
	% 4. LISTAS SEM TÍTULO AUTOMÁTICO E FORMATAÇÃO
	%------------------------------------------------------------
	\renewcommand{\listoffigures}{\@starttoc{lof}}
	\renewcommand{\listoftables}{\@starttoc{lot}}
	\renewcommand{\listof}[2]{\@starttoc{\csname ext@#1\endcsname}}
	
	\newcommand{\PreTextualTitulo}[1]{%
		\begin{center}
			\bfseries\MakeUppercase{#1}
		\end{center}
	}
	
	\makeatother
	
	% =======================================================
	% 5. IMPRESSÃO DAS LISTAS E SUMÁRIO
	% =======================================================
	
	% LISTA DE ILUSTRAÇÕES
	\SeTemItem{lof}{%
		\PreTextualTitulo{Lista de Ilustrações}
		\listoffigures
		\thispagestyle{empty}
		\newpage
	}
	
	% LISTA DE TABELAS
	\SeTemItem{lot}{%
		\PreTextualTitulo{Lista de Tabelas}
		\listoftables
		\thispagestyle{empty}
		\newpage
	}
	
	% LISTA DE ABREVIATURAS E SIGLAS
	\SeTemItem{acr}{%
		% Trava de segurança extra para garantir que não haja registros durante a impressão
		\global\def\temitem@acr@written{1}
		\PreTextualTitulo{Lista de Abreviaturas e Siglas}
		\printglossary[
		type=\acronymtype,
		style=siglasABNT,
		title={}
		]
		\thispagestyle{empty}
		\newpage
	}
	
	% LISTA DE SÍMBOLOS
	\SeTemItem{nls}{%
		\global\def\temitem@nls@written{1}
		\printnomenclature
		\thispagestyle{empty}
		\newpage
	}
	
	% LISTA DE QUADROS
	\SeTemItem{quadro}{%
		\PreTextualTitulo{Lista de Quadros}
		\listof{quadro}{}
		\thispagestyle{empty}
		\newpage
	}
	
	% LISTA DE GRÁFICOS
	\SeTemItem{grafico}{%
		\PreTextualTitulo{Lista de Gráficos}
		\listof{grafico}{}
		\thispagestyle{empty}
		\newpage
	}
	
	% SUMÁRIO
	\tableofcontents
	\thispagestyle{empty}
	
	% =======================================================
	% 6. CORREÇÃO DOS CONTADORES (ABNT)
	% =======================================================
	\setcounter{table}{0}
	\setcounter{figure}{0}
	\makeatletter
	\@ifundefined{c@quadro}{}{\setcounter{quadro}{0}}
	\@ifundefined{c@grafico}{}{\setcounter{grafico}{0}}
	\makeatother